\providecommand{\rank}{\operatorname{r}}
\providecommand{\RR}{\mathbb R}
\providecommand{\CC}{\mathbb C}
\providecommand{\KK}{\mathbb K}

% ---------- Exercise 33 ----------
\setcounter{enumi}{32}
\item {\bfseries\boldmath
State, with justification, whether each of the following assertions is true or false for any linear system $(S)$ with $m$ equations and $n$ unknowns.
}

\begin{enumerate}[label=\alph*)]

\item If $m>n$, then $(S)$ is inconsistent.

\begin{answer}
False. The system $x=1$, $2x=2$ has $m=2>1=n$ and a unique solution.
\end{answer}

\item If $m<n$, then $(S)$ is inconsistent.

\begin{answer}
False. The equation $x+y=0$ has $m=1<2=n$ and has solutions.
\end{answer}

\item If $m<n$, then $(S)$ is either inconsistent or has infinitely many solutions.

\begin{answer}
True. If the system is consistent, its number of pivots is at most
$m<n$, so at least one unknown is free. Thus it has infinitely many solutions.
Otherwise it is inconsistent.
\end{answer}

\item If $(S)$ has a unique solution, then $m=n$.

\begin{answer}
False. The counterexample in (a) has a unique solution with $m\ne n$.
\end{answer}

\item If $m=n$, then $(S)$ has a unique solution.

\begin{answer}
False. The system $x+y=0$, $2x+2y=0$ has $m=n=2$ and infinitely many
solutions. Also, $x+y=0$, $x+y=1$ is inconsistent with $m=n=2$.
\end{answer}
\end{enumerate}
